\documentclass[11pt]{article}

% ------------------------------------------------
% Packages
% ------------------------------------------------
\usepackage{amsmath,amssymb}
\usepackage{graphicx}
\usepackage{hyperref}
\usepackage{booktabs}
\usepackage{listings}
\usepackage{geometry}

\geometry{margin=1in}

% ------------------------------------------------
% Title Info
% ------------------------------------------------
\title{GraphLab: Graph Representation and Shortest Path Algorithms}
\author{Your Name}
\date{\today}

% ------------------------------------------------
\begin{document}

\maketitle

\begin{abstract}
This report documents the design and implementation of GraphLab, a small
Python framework for experimenting with graph data structures and algorithms.
We describe the chosen graph representation, analyze its complexity, and
implement fundamental algorithms including Breadth First Search and Dijkstra's
algorithm for shortest paths. The goal is to connect theoretical algorithmic
guarantees with practical implementation considerations.
\end{abstract}

% ------------------------------------------------
\section{Introduction}

Graphs are fundamental structures for modeling relationships between objects.
Applications range from networking and routing to machine learning and
optimization. Efficient graph algorithms require careful choices in data
representation and algorithm design.

The goal of GraphLab is to build a lightweight Python framework that supports

\begin{itemize}
\item graph construction
\item traversal algorithms
\item shortest path computation
\item experimentation with algorithmic performance
\end{itemize}

This report focuses on the implementation of core algorithms and their
computational properties.

% ------------------------------------------------
\section{Graph Representation}

\subsection{Adjacency List Structure}

The graph is represented using an adjacency list structure. Formally, for a
graph $G = (V,E)$, the adjacency list stores for each vertex $u \in V$ a list
of neighbors

\[
Adj[u] = \{v : (u,v) \in E\}.
\]

In the Python implementation, the adjacency list is implemented using a
dictionary mapping vertices to their neighbors.

\begin{lstlisting}[language=Python]
self.adj = {
    'A': {'B': 3, 'C': 1},
    'B': {'D': 2},
    'C': {}
}
\end{lstlisting}

This representation allows constant-time access to the neighbors of any
vertex.

\subsection{Directed and Weighted Edges}

Edges may optionally carry weights. A weighted edge is stored as

\[
Adj[u][v] = w
\]

where $w$ represents the edge weight.

% ------------------------------------------------
\section{Breadth First Search}

Breadth First Search (BFS) explores vertices in increasing distance from a
source node.

\subsection{Algorithm Description}

Starting from a source vertex $s$, BFS uses a queue to explore vertices layer
by layer.

\begin{itemize}
\item Initialize the queue with $s$
\item Mark $s$ as visited
\item Repeatedly dequeue vertices and explore their neighbors
\end{itemize}

\subsection{Correctness Intuition}

BFS discovers vertices in order of increasing distance from the source.
Therefore the first time a vertex is visited corresponds to the shortest path
in an unweighted graph.

\subsection{Complexity}

For a graph with $n = |V|$ vertices and $m = |E|$ edges:

\[
\text{BFS runtime} = O(n + m).
\]

Each vertex is processed once and each edge is examined once.

% ------------------------------------------------
\section{Shortest Paths in Unweighted Graphs}

BFS naturally computes shortest paths in unweighted graphs.

\subsection{Parent Pointer Structure}

During BFS, each vertex stores a parent pointer

\[
parent[v]
\]

representing the previous vertex in the shortest path tree.

\subsection{Path Reconstruction}

The shortest path from $s$ to a target vertex $t$ can be reconstructed by
traversing the parent pointers backwards:

\[
t \rightarrow parent[t] \rightarrow parent[parent[t]] \rightarrow \dots
\]

until the source vertex is reached.

% ------------------------------------------------
\section{Dijkstra's Algorithm}

For weighted graphs with nonnegative edge weights, shortest paths are computed
using Dijkstra's algorithm.

\subsection{Algorithm Idea}

Dijkstra's algorithm maintains a priority queue of vertices ordered by their
current distance estimate.

At each step:

\begin{enumerate}
\item Extract the vertex with minimum distance
\item Relax its outgoing edges
\item Update distances in the priority queue
\end{enumerate}

\subsection{Relaxation Step}

For an edge $(u,v)$ with weight $w$,

\[
dist[v] = \min(dist[v], dist[u] + w).
\]

\subsection{Complexity}

Using a binary heap priority queue:

\[
O(m \log n)
\]

where

\begin{itemize}
\item $n$ = number of vertices
\item $m$ = number of edges
\end{itemize}

% ------------------------------------------------
\section{Experimental Observations}

To validate the implementation, several small experiments were performed.

\subsection{Test Graphs}

Example graphs used for testing include

\begin{itemize}
\item small hand-constructed graphs
\item random graphs with varying density
\end{itemize}

\subsection{Observed Behavior}

Insert observations here once experiments are run.

Possible topics:

\begin{itemize}
\item BFS runtime scaling
\item Dijkstra performance on dense vs sparse graphs
\item memory usage
\end{itemize}

% ------------------------------------------------
\section{Discussion}

This project illustrates how theoretical algorithmic results translate into
practical implementations.

Key observations include:

\begin{itemize}
\item adjacency lists provide efficient sparse graph representation
\item BFS naturally produces shortest paths in unweighted graphs
\item Dijkstra generalizes this idea to weighted graphs
\end{itemize}

Future extensions of GraphLab could include

\begin{itemize}
\item A* search
\item minimum spanning tree algorithms
\item graph visualization tools
\end{itemize}

% ------------------------------------------------
\section{Conclusion}

GraphLab provides a simple environment for experimenting with fundamental graph
algorithms. The project highlights the relationship between theoretical
analysis and practical implementation choices.

Further work will explore additional graph algorithms and larger-scale
experiments.

% ------------------------------------------------
\end{document}